\section{Portability}
\label{sec:discussion}

\sys is designed to be portable across host and device platforms. We currently focus on the NVIDIA stack due to its widespread deployment in production, as the NVIDIA production drivers largely bypass standard Linux paths. Concrete portability, however, remains future work. On the host side, our design directly aligns with standard Linux kernel abstractions such as HMM (used by AMD ROCm with XNACK retryable faults) and the generic \texttt{migrate_vma} interface for memory management, as well as the DRM scheduler's \texttt{drm_sched_entity} abstraction for scheduling. This naturally maps onto vendor-specific scheduling units such as NVIDIA TSGs and AMD user-mode queues, whose priority and scheduling mode (e.g., \texttt{HWS} vs \texttt{HWS_NO_OVERSUBSCRIPTION}) are respected by the MES firmware during HQD mapping~\cite{rocm-gpu-memory,linux-gpusvm-doc,drm-gpu-sched,amdgpu-userq-doc}. On the device side, our compiler already implements an eBPF-to-SPIR-V backend targeting generic runtimes like AMD ROCm and Intel Level Zero. However, full end-to-end portability remains constrained by specific runtime environments, as adapting the control plane, instrumentation hooks, and performance optimizations to these platforms' particular driver interfaces and hardware characteristics is ongoing future work~\cite{spirv-ir-rfc,llvm-offload-design}.